\documentclass{beamer}
% =====================================================================
% finance-beamer.tex — shared Beamer style for the LLM-in-Finance course
% =====================================================================
% Reusable slide style. Each deck \input's this immediately after
% \documentclass{beamer}. It provides:
%   * the Madrid theme + book colour palette (matches the printed book)
%   * two book-style framing blocks:
%       \begin{bigpicture} ... \end{bigpicture}   ("the bigger picture")
%       \begin{underhood}   ... \end{underhood}    ("under the hood")
%     These mirror the book's `context` and `deepdive` tcolorboxes so the
%     same intuition-vs-mechanism scaffold the reader meets in the book
%     reappears on the slides.
%   * a Python listings style for code frames
%   * appendix conventions: \appendixstart drops a divider and switches
%     to non-numbered backup frames so "extra / less central" material
%     lives after the main talk without inflating the slide count.
%
% Usage:
%   \documentclass{beamer}
%   % =====================================================================
% finance-beamer.tex — shared Beamer style for the LLM-in-Finance course
% =====================================================================
% Reusable slide style. Each deck \input's this immediately after
% \documentclass{beamer}. It provides:
%   * the Madrid theme + book colour palette (matches the printed book)
%   * two book-style framing blocks:
%       \begin{bigpicture} ... \end{bigpicture}   ("the bigger picture")
%       \begin{underhood}   ... \end{underhood}    ("under the hood")
%     These mirror the book's `context` and `deepdive` tcolorboxes so the
%     same intuition-vs-mechanism scaffold the reader meets in the book
%     reappears on the slides.
%   * a Python listings style for code frames
%   * appendix conventions: \appendixstart drops a divider and switches
%     to non-numbered backup frames so "extra / less central" material
%     lives after the main talk without inflating the slide count.
%
% Usage:
%   \documentclass{beamer}
%   % =====================================================================
% finance-beamer.tex — shared Beamer style for the LLM-in-Finance course
% =====================================================================
% Reusable slide style. Each deck \input's this immediately after
% \documentclass{beamer}. It provides:
%   * the Madrid theme + book colour palette (matches the printed book)
%   * two book-style framing blocks:
%       \begin{bigpicture} ... \end{bigpicture}   ("the bigger picture")
%       \begin{underhood}   ... \end{underhood}    ("under the hood")
%     These mirror the book's `context` and `deepdive` tcolorboxes so the
%     same intuition-vs-mechanism scaffold the reader meets in the book
%     reappears on the slides.
%   * a Python listings style for code frames
%   * appendix conventions: \appendixstart drops a divider and switches
%     to non-numbered backup frames so "extra / less central" material
%     lives after the main talk without inflating the slide count.
%
% Usage:
%   \documentclass{beamer}
%   \input{../_style/finance-beamer.tex}
%   \title{...}\subtitle{...}\author{...}\institute{...}\date{\today}
%   \begin{document} ... \end{document}
% =====================================================================

\usetheme{Madrid}
\usecolortheme{whale}
\setbeamertemplate{navigation symbols}{}
\setbeamertemplate{caption}[numbered]
\setbeamercovered{transparent}

% --- Density controls: keep dense, equation-heavy frames on one slide ---
% Slightly smaller body text + tighter lists/blocks. Frame titles and the
% Madrid chrome keep their normal size, so the deck still reads as Madrid,
% just with more vertical room for content.
\setbeamerfont{normal text}{size=\small}
\setbeamertemplate{itemize/enumerate body begin}{\small}
\setbeamertemplate{itemize/enumerate subbody begin}{\footnotesize}
% Tighter block spacing.
\setbeamertemplate{block begin}{%
  \par\vskip2pt%
  \begin{beamercolorbox}[colsep*=.75ex,leftskip=1ex,rightskip=1ex]{block title}%
    \usebeamerfont*{block title}\insertblocktitle%
  \end{beamercolorbox}%
  {\parskip0pt\vskip-1pt}%
  \begin{beamercolorbox}[colsep*=.75ex,vmode,leftskip=1ex,rightskip=1ex]{block body}%
  \vskip1pt}
\setbeamertemplate{block end}{\end{beamercolorbox}\vskip2pt}

\usepackage{amsmath, amssymb}
\usepackage{booktabs}
\usepackage{xcolor}
\usepackage{listings}
\usepackage[skins, breakable]{tcolorbox}

% --- Book colour palette (kept in sync with book/preamble.tex) ---
\definecolor{linkblue}{RGB}{14, 75, 140}
\definecolor{deepdivebg}{RGB}{235, 244, 255}
\definecolor{narrativebg}{RGB}{255, 252, 240}
\definecolor{narrativerule}{RGB}{180, 130, 40}
\definecolor{steelblue}{RGB}{70, 130, 180}
\definecolor{forestgreen}{RGB}{34, 139, 34}
\definecolor{crimson}{RGB}{220, 20, 60}
\definecolor{darkorange}{RGB}{255, 140, 0}
\definecolor{codebg}{RGB}{248, 248, 248}
\definecolor{coderule}{RGB}{210, 210, 210}
\definecolor{keywordblue}{RGB}{0, 100, 180}
\definecolor{commentgray}{RGB}{120, 120, 120}
\definecolor{stringgreen}{RGB}{30, 130, 30}

% Tie the Madrid structural colour to the book's link blue.
\setbeamercolor{structure}{fg=linkblue}
\setbeamercolor{title}{fg=white}
\setbeamercolor{frametitle}{fg=white}

% --- "the bigger picture" : narrative / intuition framing -----------
\newtcolorbox{bigpicture}{%
  enhanced, breakable, boxsep=2pt,
  colback  = narrativebg,
  colframe = narrativerule,
  boxrule  = 0pt, toprule = 1.4pt, bottomrule = 0.4pt,
  leftrule = 0pt, rightrule = 0pt, arc = 0pt,
  left = 6pt, right = 6pt, top = 4pt, bottom = 5pt,
  fonttitle = \footnotesize\scshape\color{narrativerule},
  title = {the bigger picture}, attach title to upper = {\par\smallskip},
}

% --- "under the hood" : the mechanism / technical detail ------------
\newtcolorbox{underhood}{%
  enhanced, breakable, boxsep=2pt,
  colback  = deepdivebg,
  colframe = linkblue,
  boxrule  = 0pt, toprule = 1.4pt, bottomrule = 0.4pt,
  leftrule = 0pt, rightrule = 0pt, arc = 0pt,
  left = 6pt, right = 6pt, top = 4pt, bottom = 5pt,
  fonttitle = \footnotesize\scshape\color{linkblue},
  title = {under the hood}, attach title to upper = {\par\smallskip},
}

% --- Python code style ---------------------------------------------
\lstset{
  basicstyle    = \ttfamily\footnotesize,
  keywordstyle  = \color{keywordblue}\bfseries,
  commentstyle  = \color{commentgray}\itshape,
  stringstyle   = \color{stringgreen},
  showstringspaces = false,
  language      = Python,
  breaklines    = true,
  backgroundcolor = \color{codebg},
  frame         = single,
  rulecolor     = \color{coderule},
  numbers       = none,
  aboveskip     = 4pt, belowskip = 4pt,
}

% --- Appendix conventions ------------------------------------------
% \appendixstart{Title} : begins the "extra material" appendix. Frames
% after it are not counted toward the main deck's frame total, keeping
% the core talk lean while preserving the deeper / optional content.
\newcommand{\appendixstart}[1]{%
  \appendix
  \setbeamertemplate{frametitle continuation}{}%
  \begin{frame}[noframenumbering]
    \begin{center}
      {\Large\scshape\color{linkblue} Appendix}\\[4pt]
      {\large #1}\\[10pt]
      {\footnotesize Extra / optional material — beyond the core lecture.}
    \end{center}
  \end{frame}
}
% Use \begin{frame}[noframenumbering]{...} for each appendix frame.

%   \title{...}\subtitle{...}\author{...}\institute{...}\date{\today}
%   \begin{document} ... \end{document}
% =====================================================================

\usetheme{Madrid}
\usecolortheme{whale}
\setbeamertemplate{navigation symbols}{}
\setbeamertemplate{caption}[numbered]
\setbeamercovered{transparent}

% --- Density controls: keep dense, equation-heavy frames on one slide ---
% Slightly smaller body text + tighter lists/blocks. Frame titles and the
% Madrid chrome keep their normal size, so the deck still reads as Madrid,
% just with more vertical room for content.
\setbeamerfont{normal text}{size=\small}
\setbeamertemplate{itemize/enumerate body begin}{\small}
\setbeamertemplate{itemize/enumerate subbody begin}{\footnotesize}
% Tighter block spacing.
\setbeamertemplate{block begin}{%
  \par\vskip2pt%
  \begin{beamercolorbox}[colsep*=.75ex,leftskip=1ex,rightskip=1ex]{block title}%
    \usebeamerfont*{block title}\insertblocktitle%
  \end{beamercolorbox}%
  {\parskip0pt\vskip-1pt}%
  \begin{beamercolorbox}[colsep*=.75ex,vmode,leftskip=1ex,rightskip=1ex]{block body}%
  \vskip1pt}
\setbeamertemplate{block end}{\end{beamercolorbox}\vskip2pt}

\usepackage{amsmath, amssymb}
\usepackage{booktabs}
\usepackage{xcolor}
\usepackage{listings}
\usepackage[skins, breakable]{tcolorbox}

% --- Book colour palette (kept in sync with book/preamble.tex) ---
\definecolor{linkblue}{RGB}{14, 75, 140}
\definecolor{deepdivebg}{RGB}{235, 244, 255}
\definecolor{narrativebg}{RGB}{255, 252, 240}
\definecolor{narrativerule}{RGB}{180, 130, 40}
\definecolor{steelblue}{RGB}{70, 130, 180}
\definecolor{forestgreen}{RGB}{34, 139, 34}
\definecolor{crimson}{RGB}{220, 20, 60}
\definecolor{darkorange}{RGB}{255, 140, 0}
\definecolor{codebg}{RGB}{248, 248, 248}
\definecolor{coderule}{RGB}{210, 210, 210}
\definecolor{keywordblue}{RGB}{0, 100, 180}
\definecolor{commentgray}{RGB}{120, 120, 120}
\definecolor{stringgreen}{RGB}{30, 130, 30}

% Tie the Madrid structural colour to the book's link blue.
\setbeamercolor{structure}{fg=linkblue}
\setbeamercolor{title}{fg=white}
\setbeamercolor{frametitle}{fg=white}

% --- "the bigger picture" : narrative / intuition framing -----------
\newtcolorbox{bigpicture}{%
  enhanced, breakable, boxsep=2pt,
  colback  = narrativebg,
  colframe = narrativerule,
  boxrule  = 0pt, toprule = 1.4pt, bottomrule = 0.4pt,
  leftrule = 0pt, rightrule = 0pt, arc = 0pt,
  left = 6pt, right = 6pt, top = 4pt, bottom = 5pt,
  fonttitle = \footnotesize\scshape\color{narrativerule},
  title = {the bigger picture}, attach title to upper = {\par\smallskip},
}

% --- "under the hood" : the mechanism / technical detail ------------
\newtcolorbox{underhood}{%
  enhanced, breakable, boxsep=2pt,
  colback  = deepdivebg,
  colframe = linkblue,
  boxrule  = 0pt, toprule = 1.4pt, bottomrule = 0.4pt,
  leftrule = 0pt, rightrule = 0pt, arc = 0pt,
  left = 6pt, right = 6pt, top = 4pt, bottom = 5pt,
  fonttitle = \footnotesize\scshape\color{linkblue},
  title = {under the hood}, attach title to upper = {\par\smallskip},
}

% --- Python code style ---------------------------------------------
\lstset{
  basicstyle    = \ttfamily\footnotesize,
  keywordstyle  = \color{keywordblue}\bfseries,
  commentstyle  = \color{commentgray}\itshape,
  stringstyle   = \color{stringgreen},
  showstringspaces = false,
  language      = Python,
  breaklines    = true,
  backgroundcolor = \color{codebg},
  frame         = single,
  rulecolor     = \color{coderule},
  numbers       = none,
  aboveskip     = 4pt, belowskip = 4pt,
}

% --- Appendix conventions ------------------------------------------
% \appendixstart{Title} : begins the "extra material" appendix. Frames
% after it are not counted toward the main deck's frame total, keeping
% the core talk lean while preserving the deeper / optional content.
\newcommand{\appendixstart}[1]{%
  \appendix
  \setbeamertemplate{frametitle continuation}{}%
  \begin{frame}[noframenumbering]
    \begin{center}
      {\Large\scshape\color{linkblue} Appendix}\\[4pt]
      {\large #1}\\[10pt]
      {\footnotesize Extra / optional material — beyond the core lecture.}
    \end{center}
  \end{frame}
}
% Use \begin{frame}[noframenumbering]{...} for each appendix frame.

%   \title{...}\subtitle{...}\author{...}\institute{...}\date{\today}
%   \begin{document} ... \end{document}
% =====================================================================

\usetheme{Madrid}
\usecolortheme{whale}
\setbeamertemplate{navigation symbols}{}
\setbeamertemplate{caption}[numbered]
\setbeamercovered{transparent}

% --- Density controls: keep dense, equation-heavy frames on one slide ---
% Slightly smaller body text + tighter lists/blocks. Frame titles and the
% Madrid chrome keep their normal size, so the deck still reads as Madrid,
% just with more vertical room for content.
\setbeamerfont{normal text}{size=\small}
\setbeamertemplate{itemize/enumerate body begin}{\small}
\setbeamertemplate{itemize/enumerate subbody begin}{\footnotesize}
% Tighter block spacing.
\setbeamertemplate{block begin}{%
  \par\vskip2pt%
  \begin{beamercolorbox}[colsep*=.75ex,leftskip=1ex,rightskip=1ex]{block title}%
    \usebeamerfont*{block title}\insertblocktitle%
  \end{beamercolorbox}%
  {\parskip0pt\vskip-1pt}%
  \begin{beamercolorbox}[colsep*=.75ex,vmode,leftskip=1ex,rightskip=1ex]{block body}%
  \vskip1pt}
\setbeamertemplate{block end}{\end{beamercolorbox}\vskip2pt}

\usepackage{amsmath, amssymb}
\usepackage{booktabs}
\usepackage{xcolor}
\usepackage{listings}
\usepackage[skins, breakable]{tcolorbox}

% --- Book colour palette (kept in sync with book/preamble.tex) ---
\definecolor{linkblue}{RGB}{14, 75, 140}
\definecolor{deepdivebg}{RGB}{235, 244, 255}
\definecolor{narrativebg}{RGB}{255, 252, 240}
\definecolor{narrativerule}{RGB}{180, 130, 40}
\definecolor{steelblue}{RGB}{70, 130, 180}
\definecolor{forestgreen}{RGB}{34, 139, 34}
\definecolor{crimson}{RGB}{220, 20, 60}
\definecolor{darkorange}{RGB}{255, 140, 0}
\definecolor{codebg}{RGB}{248, 248, 248}
\definecolor{coderule}{RGB}{210, 210, 210}
\definecolor{keywordblue}{RGB}{0, 100, 180}
\definecolor{commentgray}{RGB}{120, 120, 120}
\definecolor{stringgreen}{RGB}{30, 130, 30}

% Tie the Madrid structural colour to the book's link blue.
\setbeamercolor{structure}{fg=linkblue}
\setbeamercolor{title}{fg=white}
\setbeamercolor{frametitle}{fg=white}

% --- "the bigger picture" : narrative / intuition framing -----------
\newtcolorbox{bigpicture}{%
  enhanced, breakable, boxsep=2pt,
  colback  = narrativebg,
  colframe = narrativerule,
  boxrule  = 0pt, toprule = 1.4pt, bottomrule = 0.4pt,
  leftrule = 0pt, rightrule = 0pt, arc = 0pt,
  left = 6pt, right = 6pt, top = 4pt, bottom = 5pt,
  fonttitle = \footnotesize\scshape\color{narrativerule},
  title = {the bigger picture}, attach title to upper = {\par\smallskip},
}

% --- "under the hood" : the mechanism / technical detail ------------
\newtcolorbox{underhood}{%
  enhanced, breakable, boxsep=2pt,
  colback  = deepdivebg,
  colframe = linkblue,
  boxrule  = 0pt, toprule = 1.4pt, bottomrule = 0.4pt,
  leftrule = 0pt, rightrule = 0pt, arc = 0pt,
  left = 6pt, right = 6pt, top = 4pt, bottom = 5pt,
  fonttitle = \footnotesize\scshape\color{linkblue},
  title = {under the hood}, attach title to upper = {\par\smallskip},
}

% --- Python code style ---------------------------------------------
\lstset{
  basicstyle    = \ttfamily\footnotesize,
  keywordstyle  = \color{keywordblue}\bfseries,
  commentstyle  = \color{commentgray}\itshape,
  stringstyle   = \color{stringgreen},
  showstringspaces = false,
  language      = Python,
  breaklines    = true,
  backgroundcolor = \color{codebg},
  frame         = single,
  rulecolor     = \color{coderule},
  numbers       = none,
  aboveskip     = 4pt, belowskip = 4pt,
}

% --- Appendix conventions ------------------------------------------
% \appendixstart{Title} : begins the "extra material" appendix. Frames
% after it are not counted toward the main deck's frame total, keeping
% the core talk lean while preserving the deeper / optional content.
\newcommand{\appendixstart}[1]{%
  \appendix
  \setbeamertemplate{frametitle continuation}{}%
  \begin{frame}[noframenumbering]
    \begin{center}
      {\Large\scshape\color{linkblue} Appendix}\\[4pt]
      {\large #1}\\[10pt]
      {\footnotesize Extra / optional material — beyond the core lecture.}
    \end{center}
  \end{frame}
}
% Use \begin{frame}[noframenumbering]{...} for each appendix frame.


% Self-contained citation fallback (no .bib pass for these decks).
\newcommand{\citesrc}[1]{\textsc{\small #1}}
\providecommand{\citet}[1]{\citesrc{#1}}
\providecommand{\citep}[1]{{\footnotesize[\,\citesrc{#1}\,]}}

\title{Practical Session 13: Evaluating LLMs Honestly}
\subtitle{Lecture 13 --- Calibration, Contamination Audits, and Trading-Grade Metrics}
\author{Juan F.\ Imbet}
\institute{Paris Dauphine -- PSL University}
\date{\today}

\begin{document}

\begin{frame}
  \titlepage
\end{frame}

% ===============================================================
\begin{frame}{Session Overview}
\textbf{Duration:} $\sim$1 hour (50 min hands-on + 10 min debrief)

\medskip
\begin{block}{Timed agenda}
\begin{enumerate}
  \item \textbf{Recap} (5 min): ECE, temporal splits, Sharpe/IR, transaction costs.
  \item \textbf{Problem 1} (12 min): Compute ECE \& MCE by hand from bin statistics.
  \item \textbf{Problem 2} (12 min): Design a contamination-resistant split; quantify
        a leakage audit.
  \item \textbf{Problem 3} (10 min): Gross vs.\ net Sharpe under transaction costs.
  \item \textbf{Case study} (15 min): Python --- reliability diagram + temperature
        scaling on a toy LLM credit classifier
        (\texttt{code/notebooks/13-llm-limitations-evaluation/}).
  \item \textbf{Discussion} (6 min): Would you deploy this model?
\end{enumerate}
\end{block}
\end{frame}

% ===============================================================
\begin{frame}{Recap: The Four Formulas You Need Today}
\textbf{Expected / Maximum Calibration Error:}
\[
  \text{ECE} = \sum_{m=1}^{M}\frac{|B_m|}{n}\bigl|\text{acc}(B_m)-\overline{p}(B_m)\bigr|,
  \quad
  \text{MCE} = \max_m \bigl|\text{acc}(B_m)-\overline{p}(B_m)\bigr|.
\]

\textbf{Temporal split with gap} $\Delta \geq h$:
\[
  \text{train}: t \leq T_1-\Delta,\quad
  \text{valid}: T_1\!\leq t \leq T_2-\Delta,\quad
  \text{test}: t \geq T_2.
\]

\textbf{Sharpe and net-of-cost return:}
\[
  \text{SR} = \frac{\overline{R}_p-R_f}{\sigma_p},\qquad
  R^{\text{net}}_p = R^{\text{gross}}_p - 2\tau(\text{half-spread}+\text{impact}).
\]

\medskip
\begin{block}{Recall}
Temperature scaling $\operatorname{softmax}(\mathbf{z}/T)$ with $T>1$ cures
overconfidence and \emph{preserves ranking} (AUROC unchanged).
\end{block}
\end{frame}

% ===============================================================
\begin{frame}{Problem 1: Compute ECE and MCE}
An LLM sentiment model labels 1{,}000 headlines as ``price-positive''. Confidence
binned into 4 bins:

\medskip
\centering
\begin{tabular}{ccrcc}
\toprule
Bin & Interval & $|B_m|$ & $\overline{p}(B_m)$ & $\text{acc}(B_m)$ \\
\midrule
1 & [0.0, 0.25) & 100 & 0.15 & 0.18 \\
2 & [0.25, 0.5) & 200 & 0.40 & 0.35 \\
3 & [0.5, 0.75) & 300 & 0.65 & 0.50 \\
4 & [0.75, 1.0] & 400 & 0.90 & 0.66 \\
\bottomrule
\end{tabular}

\medskip
\raggedright
\textbf{Tasks.}
\begin{enumerate}
  \item Compute the per-bin gap $|\text{acc}-\overline{p}|$.
  \item Compute \textbf{ECE} (weight by $|B_m|/n$) and \textbf{MCE}.
  \item In which confidence region is the model worst? What does that imply for a
        desk that only acts on $>0.75$ signals?
\end{enumerate}
\end{frame}

% ===============================================================
\begin{frame}{Solution: Problem 1}
\textbf{Per-bin gaps:} $0.03,\ 0.05,\ 0.15,\ 0.24$. \quad
\textbf{Weights:} $0.10,\ 0.20,\ 0.30,\ 0.40$.

\medskip
\[
  \text{ECE} = 0.10(0.03)+0.20(0.05)+0.30(0.15)+0.40(0.24)
\]
\[
  = 0.003 + 0.010 + 0.045 + 0.096 = \mathbf{0.154}.
\]
\[
  \text{MCE} = \max\{0.03,0.05,0.15,0.24\} = \mathbf{0.24}.
\]

\medskip
\begin{alertblock}{Interpretation}
ECE $=15.4\%$, MCE $=24\%$, both in the top bin. The desk acting only on
$>0.75$ signals is operating in the model's \textbf{worst} region: stated 90\%
confidence corresponds to \textbf{66\%} realised accuracy. Recalibrate (or
abstain) before trusting high-confidence calls.
\end{alertblock}
\end{frame}

% ===============================================================
\begin{frame}{Problem 2: Design a Contamination Audit}
A team evaluates an LLM (training cutoff \textbf{Dec 2023}) on predicting 1-month
stock direction from earnings calls. Prediction horizon $h = 21$ trading days.

\medskip
\textbf{Tasks.}
\begin{enumerate}
  \item Give concrete cutoffs $T_1, T_2$ and gap $\Delta$ for a temporal split that
        is safe against this model. Justify $\Delta$.
  \item The team measures accuracy on \emph{original} (62\%) vs.\
        \emph{anonymised} (54\%) transcripts. Quantify the leakage and state what
        it implies about the true edge.
  \item Name two further controls and what each detects.
\end{enumerate}
\end{frame}

% ===============================================================
\begin{frame}{Solution: Problem 2}
\textbf{(1) Split.} All test data must post-date the cutoff plus a publication-lag
margin. Use $T_2 = $ \textbf{Apr 2024} (Dec 2023 cutoff $+ \geq 3$ months);
$T_1$ in early 2023 for validation. $\Delta \geq h + \text{lag} \approx
21\text{d} + \text{a few weeks} \Rightarrow$ take $\Delta = 1$ month so labels
cannot straddle the boundary.

\medskip
\textbf{(2) Leakage.} $62\% - 54\% = \mathbf{8}$ pp of accuracy was
entity-level memorisation. Residual 54\% still beats 50\% but the \emph{apparent}
edge shrank by more than half---likely economically negligible after costs.

\medskip
\textbf{(3) Two further controls.}
\begin{itemize}
  \item \textbf{N-gram dedup} (13-grams, MinHash LSH) --- detects direct
        memorisation of the document text.
  \item \textbf{Date-blinding ablation} --- a large drop when temporal markers are
        removed reveals temporal pattern-retrieval rather than reasoning.
\end{itemize}
\end{frame}

% ===============================================================
\begin{frame}{Problem 3: Gross vs.\ Net Sharpe}
A daily news-sentiment LLM strategy reports, before costs:
\[
  \overline{R}^{\text{gross}}_p - R_f = 12\%\ \text{annual},\qquad
  \sigma_p = 15\%\ \text{annual}.
\]
It rebalances daily with turnover $\tau = 0.5$/day; half-spread $= 2$ bps,
market impact $= 8$ bps.

\medskip
\textbf{Tasks.}
\begin{enumerate}
  \item Compute the \textbf{gross} annual Sharpe ratio.
  \item Daily round-trip cost $= 2\tau(\text{half-spread}+\text{impact})$.
        Annualise over 252 trading days.
  \item Compute the \textbf{net} Sharpe (subtract annual cost from excess return;
        $\sigma_p$ unchanged). Is the strategy still worth trading?
\end{enumerate}
\end{frame}

% ===============================================================
\begin{frame}{Solution: Problem 3}
\textbf{(1) Gross Sharpe:} $\text{SR}_{\text{gross}} = 12/15 = \mathbf{0.80}$.

\medskip
\textbf{(2) Cost.} Daily round-trip $= 2(0.5)(2+8)\ \text{bps} = 10$ bps/day.
Over 252 days: $10\ \text{bps}\times 252 = 2520\ \text{bps} = \mathbf{25.2\%}$/year.

\medskip
\textbf{(3) Net Sharpe.}
\[
  \overline{R}^{\text{net}}_p - R_f = 12\% - 25.2\% = -13.2\%,
  \qquad
  \text{SR}_{\text{net}} = \frac{-13.2}{15} = \mathbf{-0.88}.
\]

\medskip
\begin{alertblock}{Verdict}
A respectable-looking gross Sharpe of 0.80 turns \textbf{deeply negative} once
daily-rebalancing costs are charged. The ``edge'' was an artefact of ignoring
transaction costs---exactly the artefact the lecture warns about. Lower turnover
or longer holding periods are mandatory before this signal is tradeable.
\end{alertblock}
\end{frame}

% ===============================================================
\begin{frame}{Case Study: Reliability + Temperature Scaling}
\textbf{Goal.} Reproduce the credit-classifier reliability diagram, compute its
ECE, then apply temperature scaling and watch ECE fall while AUROC is unchanged.

\medskip
\textbf{Reference code:}
\texttt{code/notebooks/13-llm-limitations-evaluation/}
\begin{itemize}
  \item \texttt{gen\_reliability\_diagram.py} --- deterministic diagram from the
        fixed bin statistics (ECE $=0.115$).
  \item \texttt{exercises.ipynb} / \texttt{demo.ipynb} --- calibration experiments.
\end{itemize}

\medskip
\begin{block}{What you will observe}
\begin{enumerate}
  \item The high-confidence bins sit \emph{below} the diagonal --- overconfidence.
  \item Fitting $T>1$ spreads the softmax: bins move toward the diagonal, ECE drops.
  \item The ROC curve (and AUROC) is \textbf{unchanged} --- scaling is
        rank-preserving. Calibration $\neq$ discrimination.
\end{enumerate}
\end{block}
\end{frame}

% ===============================================================
\begin{frame}[fragile]{Case Study: The Code}
\begin{lstlisting}
import numpy as np

# Fixed bin stats from the credit example (ex:ch13-ece-credit).
conf  = np.array([0.08, 0.31, 0.50, 0.71, 0.88])  # mean predicted prob
acc   = np.array([0.11, 0.28, 0.47, 0.55, 0.62])  # empirical accuracy
count = np.array([120, 180, 200, 300, 200])

def ece(conf, acc, count):
    w = count / count.sum()
    return float(np.sum(w * np.abs(conf - acc)))

print("ECE  =", round(ece(conf, acc, count), 3))   # -> 0.115
print("MCE  =", round(float(np.max(np.abs(conf - acc))), 3))

# Temperature scaling on raw logits z: p = softmax(z / T).
# T > 1 spreads the distribution and reduces overconfidence;
# because it is monotone in z, the ranking (AUROC) is preserved.
def temperature_scale(logits, T):
    z = logits / T
    z = z - z.max(axis=1, keepdims=True)        # numerical stability
    e = np.exp(z)
    return e / e.sum(axis=1, keepdims=True)
\end{lstlisting}
\end{frame}

% ===============================================================
\begin{frame}{Case Study: Diagnostic Questions}
Run \texttt{gen\_reliability\_diagram.py}, then answer:

\medskip
\textbf{Q1.} Which bins lie \emph{below} the diagonal, and what does that say about
the model's confidence in that region?

\medskip
\textbf{Q2.} The reported ECE is $0.115$. Recompute it using only bins 4 and 5.
Why do these two bins dominate the total?

\medskip
\textbf{Q3.} You fit $T = 1.6$ and ECE falls to $\approx 0.04$, but AUROC is
unchanged at $0.78$. Explain in one sentence why calibration improved while
discrimination did not.

\medskip
\textbf{Q4.} A colleague proposes raising the action threshold from $0.5$ to
$0.8$ instead of recalibrating. Does that fix the calibration problem? What does
it fix, and what does it not?
\end{frame}

% ===============================================================
\begin{frame}{Discussion: Would You Deploy This Model?}
You are on a model-risk committee (SR~11-7). An LLM credit model shows:
ECE $=0.115$, 8 pp of its accuracy attributable to entity memorisation, and a
news-overlay trading sleeve with gross Sharpe 0.80 / net Sharpe $-0.88$.

\medskip
\begin{block}{Debate, framed around}
\begin{itemize}
  \item \textbf{Calibration} --- deploy with recalibration + conformal sets, or not?
  \item \textbf{Contamination} --- is the residual 54\% edge real or an artefact?
  \item \textbf{Economics} --- can lower turnover rescue the trading sleeve?
  \item \textbf{Governance} --- what monitoring cadence (quarterly recalibration?)
        and human-in-the-loop controls would you require before sign-off?
\end{itemize}
\end{block}

\medskip
\textbf{Takeaway:} LLM outputs are \emph{inputs to a decision process}, not end
products. The committee's job is the surrounding evaluation discipline.
\end{frame}

% ===============================================================
\appendixstart{Stretch Problems}
% ===============================================================

\begin{frame}[noframenumbering]{Stretch: Information Ratio Decomposition}
A long-short LLM signal covers $N = 100$ independent monthly bets with
Information Coefficient $\text{IC} = 0.04$.

\medskip
\textbf{Tasks.}
\begin{enumerate}
  \item Estimate the Information Ratio via $\text{IR} \approx \text{IC}\sqrt{N}$.
  \item Is this ``strong'' (IR $>0.5$) or ``exceptional'' (IR $>1.0$)?
  \item To reach IR $=0.8$ at the same IC, how many bets $N$ would you need? What
        practical cost does increasing $N$ introduce?
\end{enumerate}

\medskip
\textbf{Solution sketch.}
$\text{IR} \approx 0.04\sqrt{100} = 0.40$ (below ``strong''). For IR $=0.8$:
$\sqrt{N} = 0.8/0.04 = 20 \Rightarrow N = 400$ bets/year. More bets means higher
turnover $\Rightarrow$ higher transaction costs, which can erase the gain
(see Problem 3).
\end{frame}

\begin{frame}[noframenumbering]{Stretch: Walk-Forward with a Purge Period}
You backtest a monthly-rebalanced LLM signal ($h = 1$ month) over 2019--2024.

\medskip
\textbf{Tasks.}
\begin{enumerate}
  \item Sketch the train/eval windows for the first three walk-forward iterations
        with purge $\Delta_{\text{purge}} = 1$ month.
  \item Why must $\Delta_{\text{purge}} \geq h$? What leaks if it is zero?
  \item For full LLM fine-tuning (expensive to re-estimate), what is the cost of
        widening $h$ to amortise re-training?
\end{enumerate}

\medskip
\textbf{Solution sketch.}
Iteration $k$: train $t \leq kh - \Delta_{\text{purge}}$, eval
$kh \leq t \leq (k{+}1)h-1$. With $\Delta_{\text{purge}} < h$, a label window
overlapping the boundary appears in both sets $\Rightarrow$ look-ahead leakage.
Widening $h$ amortises re-fitting cost but evaluates an increasingly \emph{stale}
model late in each window.
\end{frame}

\end{document}
